\section{The gradient atlas}
\label{app:atlas}

This appendix carries the detail behind \S\ref{sec:eikonal}: the measurement that
told us the shipped shader was omitting the eikonal divide, and how wrong that was
per family.

\subsection{What was measured}

For a family given as the level sets of $\ph$, the Euclidean distance to the
nearest member is $\mathrm{pd}(\ph-\varphi,s)/\lVert\nabla\ph\rVert$. Drop the
denominator and the value you threshold is a phase residual, not a length, so the
stroke is too thin by exactly $\lVert\nabla\ph\rVert$ wherever that exceeds $1$.
The measurement is therefore direct: sweep each curve family over its own
parameter ranges, evaluate $\lVert\nabla\ph\rVert$ on a $1400\times880$ frame, and
report its worst value and the share of the frame where it departs from $1$ by
more than $10\%$.

% GENERATED by paper/tools/numbers.mjs
\begin{table}[t]
\caption{\figlead{How wrong an undivided phase residual is.} For each curve
family, whether the shipped shader divided by $\lVert\nabla\ph\rVert$, how many
settings of that family's own sliders we swept, the worst
$\lVert\nabla\ph\rVert$ found on a $1400\times880$ frame---which is exactly
the factor by which the stroke was too thin---and the share of the frame where
the stroke was off by more than $10\%$. The parabola was correct all along; the
other three were not.}
\label{tab:gradient}
\small
\begin{tabular}{@{}lrrrr@{}}
\toprule
Family & divides & settings & worst factor & frame off $>10\%$ \\
\midrule
Wave & \textbf{no} & 36 & 25.2$\times$ & 61\% \\
Parabola & yes & 6 & 224.0$\times$ & 98\% \\
Hyperbola & \textbf{no} & 1 & 14.8$\times$ & 66\% \\
Spiral & \textbf{no} & 5 & 7.3$\times$ & 0\% \\
\bottomrule
\end{tabular}
\end{table}


Table~\ref{tab:gradient} is the result. Three of the four curve families divided
by nothing, and the worst factors are not marginal: a wave at the top of its
amplitude range is $\gradWorstWave\times$ too thin, a hyperbola
$\gradWorstHyp\times$, a spiral $\gradWorstSpiral\times$.

\subsection{Reading the table the other way}

Two rows repay a second look, and between them they say why this bug survived so
long.

\paragraph{The spiral was wrong and looked fine.} Its worst factor is
$\gradWorstSpiral\times$, but the share of the frame off by more than $10\%$ is
$\gradSpiralFrame\%$. For an Archimedean spiral $\lVert\nabla\ph\rVert=\sqrt{1+b^2/r^2}$, which is
close to $1$ everywhere except very near the origin --- and near the origin the
arms are so tightly packed that nobody was looking. The bug was real, confined to
a region already illegible, and therefore invisible.

\paragraph{The parabola was right and looks worst.} It divides, and always did,
yet it has the largest worst factor in the table at $\gradWorstParab\times$ and
$\gradParabFrame\%$ of the frame off. Those two numbers describe the family, not the
bug: a parabola's members genuinely crowd towards the edges of a wide frame, so
$\lVert\nabla\ph\rVert$ genuinely reaches $\gradWorstParab$, and the correct
renderer genuinely divides by it. Had we
read only the ``worst factor'' column we would have concluded the parabola was the
most broken family in the tool. It is the only one that was never broken.

The lesson we take is about the shape of the measurement rather than the numbers.
A per-family scalar --- worst case, mean, percentile --- cannot separate ``this
family is steep'' from ``this renderer is wrong,'' because the same quantity
appears in both. What separates them is comparing the two renderers on the same
family, which is what Figure~\ref{fig:gradient} does and what we should have done
first.

\subsection{Consequence for the stroke floor}

The floor of \eqref{eq:stroke} is $1.15$ pixels, in world units. Applied to a
phase residual it means $1.15$ phase units, which is $1.15/\lVert\nabla\ph\rVert$
world units --- so on a steep family the protection against zoom-out speckle was
weaker than intended by the same factor the stroke was thin by, and weakest
exactly where the family was densest and needed it most. The two bugs are one bug,
and fixing the divide fixed both. We mention it because the failure mode of the
floor is the same visual signature as the inversion failure of
\S\ref{sec:inverse} --- a pattern that breaks into speckle at low zoom --- and we
spent real time attributing one to the other.

\subsection{The same divide, for an authored field}

A modulated family (\S\ref{sec:beyond}) has phase $\ph - a s f$, so its gradient is
$\nabla\ph - a s \nabla f$ and the divide needs $\nabla f$ at every pixel. This is
the reason \S\ref{sec:exprs} differentiates rather than asking the author to: a
field supplied without its gradient reintroduces exactly the bug this appendix is
about, and does so on a term the author controls with a slider. The presets are
listed as source in Table~\ref{tab:exprs}; the gradients are not listed anywhere,
because there is nothing to list --- both evaluators carry the partials alongside
the value and neither has a derivative written down in it.

% GENERATED by paper/tools/numbers.mjs
\begin{table}[t]
  \centering
  \small
  \caption{The presets the field editor starts from, as the text a user types.
  Nothing here is privileged: each is parsed, compiled and unrolled down the same
  path any typed expression takes, and each is numerically the hand-written field
  it replaced rather than a lookalike.}
  \label{tab:exprs}
  \begin{tabular}{@{}lp{0.66\linewidth}@{}}
    \toprule
    Preset & Expression \\
    \midrule
    \textsc{saddle} & {\ttfamily\scriptsize x\textasciicircum{}2 - \allowbreak y\textasciicircum{}2} \\
    \textsc{dipole} & {\ttfamily\scriptsize 0.5 * (1 / sqrt((x - \allowbreak 1)\textasciicircum{}2 + \allowbreak y\textasciicircum{}2 + \allowbreak 0.0625) - \allowbreak 1 / sqrt((x + \allowbreak 1)\textasciicircum{}2 + \allowbreak y\textasciicircum{}2 + \allowbreak 0.0625))} \\
    \textsc{bumps} & {\ttfamily\scriptsize exp(-((x + \allowbreak 0.45)\textasciicircum{}2 + \allowbreak (y - \allowbreak 0.35)\textasciicircum{}2) / (2 * 0.45\textasciicircum{}2)) - \allowbreak 1.2 * exp(-((x - \allowbreak 0.55)\textasciicircum{}2 + \allowbreak (y + \allowbreak 0.4)\textasciicircum{}2) / (2 * 0.55\textasciicircum{}2)) + \allowbreak 0.65 * exp(-((x - \allowbreak 0.15)\textasciicircum{}2 + \allowbreak (y - \allowbreak 0.78)\textasciicircum{}2) / (2 * 0.33\textasciicircum{}2))} \\
    \textsc{swirl} & {\ttfamily\scriptsize -0.5 * log((x + \allowbreak 0.6)\textasciicircum{}2 + \allowbreak (y + \allowbreak 0.7)\textasciicircum{}2 + \allowbreak 0.0324) + \allowbreak 0.5 * log((x - \allowbreak 0.62)\textasciicircum{}2 + \allowbreak (y + \allowbreak 0.66)\textasciicircum{}2 + \allowbreak 0.0324) - \allowbreak 0.425 * log((x - \allowbreak 0.55)\textasciicircum{}2 + \allowbreak (y - \allowbreak 0.7)\textasciicircum{}2 + \allowbreak 0.0324) + \allowbreak 0.35 * log((x + \allowbreak 0.58)\textasciicircum{}2 + \allowbreak (y - \allowbreak 0.72)\textasciicircum{}2 + \allowbreak 0.0324)} \\
    \textsc{ripple} & {\ttfamily\scriptsize cos(tau * r)} \\
    \textsc{terrain} & {\ttfamily\scriptsize 0.34 * sin(1.9 * x + \allowbreak 1.3 * y + \allowbreak 0.3) + \allowbreak 0.24 * sin(3.1 * x - \allowbreak 2.4 * y + \allowbreak 1.9) + \allowbreak 0.15 * sin(5 * x + \allowbreak 3.9 * y + \allowbreak 3.4) + \allowbreak 0.13 * sin(-1.3 * x + \allowbreak 5.8 * y + \allowbreak 5.1) + \allowbreak 0.09 * sin(7.4 * x - \allowbreak 1.35 * y + \allowbreak 2.2)} \\
    \bottomrule
  \end{tabular}
\end{table}

